\documentclass[12pt, letterpaper, parskip = full]{article}
\usepackage[mmddyyyy]{datetime}

\usepackage[utf8]{inputenc}
\usepackage{amsmath}
\usepackage{mathpazo}
\usepackage{eulervm}
\usepackage[skip=10pt plus1pt]{parskip}
\usepackage[normalem]{ulem}

\usepackage[margin=1in]{geometry}
\linespread{1.20}
\usepackage{enumerate}
\usepackage{enumitem}
\usepackage{listings}
\usepackage[T1]{fontenc}
\usepackage{amsfonts}
\usepackage{amssymb}
\usepackage{amsthm}
\usepackage{multirow,array}
\usepackage{graphicx}
\usepackage{physics}
\usepackage{bm}
\usepackage{tikz}
\usepackage{breqn}
\usepackage[longnamesfirst]{natbib}
\usepackage{comment}
\usepackage{setspace}
\usepackage{hyperref}
\usepackage{accents}
\usepackage{graphicx}
\usepackage{tabularx}

\graphicspath{{Images/}}
\setlength{\parindent}{0pt}

\newtheorem{theorem}{Theorem}
\newtheorem{corollary}{Corollary}
\newtheorem{lemma}{Lemma}
\newtheorem{definition}{Definition}
\newtheorem{axiom}{Axiom}
\newcommand{\e}{\mathrm{e}}
\newcommand{\Var}{\mathrm{Var}}
\newcommand{\Ex}[1]{\mathbb{E}\left[#1\right]}
\renewcommand{\Pr}[1]{\mathbb{P}(#1)}
\newcommand{\plim}{\mathrm{plim}}
\newcommand{\pc}[1]{\frac{\dot{#1}(t)}{#1(t)}}
\renewcommand\qedsymbol{$\blacksquare$}
\newcommand{\ubar}[1]{\underaccent{\bar}{#1}}

\title{Homework 10}
\author{Alejandro Flores-Brown}

\begin{document}

	\maketitle

    \section{H 10-1}

    This is the same project from the worksheet, so we have a \(50\%\) chance of a cashflow of \(140\) and a \(50\%\) chance of a cashflow of \(90\), which means that after discounting at a weighted average cost of capital of \(15\%\) our cashflow has a present value of \(\frac{0.5 \times (140 + 90)}{1.15} = 100\). We assume the entrepreneur raises \(70\) in debt, so the equityholder's cashflow is \( 0.5 \times 0.7 \times (140 + 90) = 80.5 \) and they need to hold \(\frac{80.5}{1.15} = 70\%\) of the company to be adequately compensated for their risk.

    \section{H 10-2}

    We are given that the debt-to-equity ratio of \(0.4\), an equity beta of \(1.4\), and a debt beta of \(0.2\). We can interpret the first as \(\frac{D}{E} = 0.4\), which means all further ratios will be satisfied if we write \(D = 0.4\) and \(E = 1\). That lets us solve for the unlevered equity beta:

    \[ \beta_U = \frac{E}{D+E} \beta_E + \frac{D}{D+E} \beta_D = \frac{1}{1.4} \times 1.4 + \frac{0.4}{1.4} \times 0.2 = 1 + \frac{0.4}{7} = \frac{1.4}{7} = 0.2   \]

    Now we hold this fixed and reverse-engineer the equity beta of \(1.2\):
    
    \[ 0.2 = \frac{1}{1.4} \times 1.2 + \frac{0.4}{1.4} \beta_D \implies \beta_D = \frac{7}{2} \left( 0.2 - \frac{5}{7} \right) = \frac{7}{10} - \frac{5}{2} = 0.7 - 2.5 = -1.8  \]



    \section{H 10-3}

    The weighted average cost of capital is given by

    \[ r_U = \frac{E}{D+E-C} r_E + \frac{D}{D+E-C} r_D - \frac{C}{D+E-C} r_f = \frac{20}{20 + 40 - 4} \left( 0.04 + 1.2 \times 0.06 \right) + \frac{40}{20 + 40 - 4} \left( 0.04 + 0.1 \times 0.06 \right) - \frac{4}{20 + 40 - 4} \times 0.04  \]

    If we instead use the cash to clear debt and hold the betas constant, we instead can find an unlevered equity beta so that we have a beta for the asset:

    \[ \beta_U = \frac{40}{40 + 16} \times 1.2 + \frac{16}{40 + 16} \times 0.1    \]


\end{document}